% GENERATED FILE. Edit canonical lessons, then run npm run generate:books.
% canonical-source-hash: fnv1a64:9474ca80916b5cc3
% canonical-lessons: JA-C201-wataru, JA-C201-komu, JA-C201-yokogiru, JA-C201-tobidasu, JA-C201-butsukaru

\chapter{Crossing the Road Safely}
\label{ch:ja-crossing-the-road-safely}
\hlchaptermodality{201}

\begin{chapteropening}
\textbf{By the end of this chapter:} \emph{I can say how people cross roads and what happens in traffic.}

Complete the last lesson of chapter 201: I can say how people cross roads and what happens in traffic.
\end{chapteropening}

\section[\texorpdfstring{\ja{わたる} (wataru)}{わたる (wataru)}]{\ja{わたる} (wataru) — to cross; to go across}

\label{lesson:JA-C201-wataru}



\begin{quote}
Before the new one: say the Japanese for a boarding point, then the Japanese for a bus stop.
\end{quote}

\subsection*{You'll want to know: \ja{わたる}}
\textbf{\ja{わたる}} — \emph{wataru} — \textquotedblleft{}to cross; to go across\textquotedblright{}.

\begin{quote}
\textbf{\ja{しんごうで} \ja{どうろを} \ja{わたります。}} — \emph{shingō de dōro o watarimasu} — \textquotedblleft{}I cross the road at the traffic lights.\textquotedblright{}
\end{quote}

Polite \textbf{\ja{わたります}}, \emph{watarimasu}; a godan verb.

\begin{cousinweb}
\textbf{\ja{わたる}} (\emph{wataru}) is a native Japanese word (wago), inherited rather than borrowed from Chinese; its kanji was attached later as a kun'yomi, and its older origin is not settled (JMdict, KANJIDIC2).
\end{cousinweb}

\begin{grammarlens}[title={What you've built}]
One more word for everyday life at A2.
\end{grammarlens}

\subsection*{Your turn}
\begin{itemize}
\raggedright
  \item \emph{Say it:} \emph{wataru}
  \item \emph{Say it:} \emph{wataru}, once more
  \item \emph{Say it:} \emph{shingō de dōro o watarimasu} — I cross the road at the traffic lights.
  \item \emph{Recall:} say the Japanese for a boarding point, then the Japanese for a bus stop, then say \emph{wataru} again
\end{itemize}

\begin{tcolorbox}[breakable,colback=teal!4,colframe=teal!35!black,title={Before you move on}]
What does \ja{わたる} mean? (\textquotedblleft{}To cross; to go across\textquotedblright{}.) Say it once more.
\end{tcolorbox}
\section[\texorpdfstring{\ja{こむ} (komu)}{こむ (komu)}]{\ja{こむ} (komu) — to be crowded}

\label{lesson:JA-C201-komu}



\begin{quote}
Before the new one: say the Japanese for a bus stop, then the Japanese for to cross.
\end{quote}

\subsection*{You'll want to know: \ja{こむ}}
\textbf{\ja{こむ}} — \emph{komu} — \textquotedblleft{}to be crowded\textquotedblright{}.

\begin{quote}
\textbf{\ja{あさの} \ja{でんしゃは} \ja{こみます。}} — \emph{asa no densha wa komimasu} — \textquotedblleft{}Morning trains get crowded.\textquotedblright{}
\end{quote}

Polite \textbf{\ja{こみます}}, \emph{komimasu}; a godan verb.

\begin{cousinweb}
\textbf{\ja{こむ}} (\emph{komu}) is a native Japanese word (wago), inherited rather than borrowed from Chinese; its kanji was attached later as a kun'yomi, and its older origin is not settled (JMdict, KANJIDIC2).
\end{cousinweb}

\begin{grammarlens}[title={What you've built}]
One more word for everyday life at A2.
\end{grammarlens}

\subsection*{Your turn}
\begin{itemize}
\raggedright
  \item \emph{Say it:} \emph{komu}
  \item \emph{Say it:} \emph{komu}, once more
  \item \emph{Say it:} \emph{asa no densha wa komimasu} — Morning trains get crowded.
  \item \emph{Recall:} say the Japanese for a bus stop, then the Japanese for to cross, then say \emph{komu} again
\end{itemize}

\begin{tcolorbox}[breakable,colback=teal!4,colframe=teal!35!black,title={Before you move on}]
What does \ja{こむ} mean? (\textquotedblleft{}To be crowded\textquotedblright{}.) Say it once more.
\end{tcolorbox}
\section[\texorpdfstring{\ja{よこぎる} (yokogiru)}{よこぎる (yokogiru)}]{\ja{よこぎる} (yokogiru) — to cut across; to cross (a road)}

\label{lesson:JA-C201-yokogiru}



\begin{quote}
Before the new one: say the Japanese for to cross, then the Japanese for to be crowded.
\end{quote}

\subsection*{You'll want to know: \ja{よこぎる}}
\textbf{\ja{よこぎる}} — \emph{yokogiru} — \textquotedblleft{}to cut across; to cross (a road)\textquotedblright{}.

\begin{quote}
\textbf{\ja{こどもが} \ja{どうろを} \ja{よこぎります。}} — \emph{kodomo ga dōro o yokogirimasu} — \textquotedblleft{}A child cuts across the road.\textquotedblright{}
\end{quote}

Polite \textbf{\ja{よこぎります}}, \emph{yokogirimasu}; a godan verb.

\begin{cousinweb}
\textbf{\ja{よこぎる}} (\emph{yokogiru}) is native Japanese (wago), built from \emph{yoko}, \textquotedblleft{}side\textquotedblright{}, and \emph{kiru}, \textquotedblleft{}to cut\textquotedblright{}, both inherited rather than borrowed; the \emph{k} is voiced to \emph{g} (JMdict).
\end{cousinweb}

\begin{grammarlens}[title={What you've built}]
One more word for everyday life at A2.
\end{grammarlens}

\subsection*{Your turn}
\begin{itemize}
\raggedright
  \item \emph{Say it:} \emph{yokogiru}
  \item \emph{Say it:} \emph{yokogiru}, once more
  \item \emph{Say it:} \emph{kodomo ga dōro o yokogirimasu} — A child cuts across the road.
  \item \emph{Recall:} say the Japanese for to cross, then the Japanese for to be crowded, then say \emph{yokogiru} again
\end{itemize}

\begin{tcolorbox}[breakable,colback=teal!4,colframe=teal!35!black,title={Before you move on}]
What does \ja{よこぎる} mean? (\textquotedblleft{}To cut across; to cross (a road)\textquotedblright{}.) Say it once more.
\end{tcolorbox}
\section[\texorpdfstring{\ja{とびだす} (tobidasu)}{とびだす (tobidasu)}]{\ja{とびだす} (tobidasu) — to rush out; to jump out}

\label{lesson:JA-C201-tobidasu}



\begin{quote}
Before the new one: say the Japanese for to be crowded, then the Japanese for to cut across.
\end{quote}

\subsection*{You'll want to know: \ja{とびだす}}
\textbf{\ja{とびだす}} — \emph{tobidasu} — \textquotedblleft{}to rush out; to jump out\textquotedblright{}.

\begin{quote}
\textbf{\ja{こどもが} \ja{みちに} \ja{とびだします。}} — \emph{kodomo ga michi ni tobidashimasu} — \textquotedblleft{}A child rushes out into the street.\textquotedblright{}
\end{quote}

Polite \textbf{\ja{とびだします}}, \emph{tobidashimasu}; a godan verb.

\begin{cousinweb}
\textbf{\ja{とびだす}} (\emph{tobidasu}) is native Japanese (wago), built from \emph{tobi}, \textquotedblleft{}flying, leaping\textquotedblright{}, and \emph{dasu}, \textquotedblleft{}to put out\textquotedblright{}, both inherited rather than borrowed (JMdict).
\end{cousinweb}

\begin{grammarlens}[title={What you've built}]
One more word for everyday life at A2.
\end{grammarlens}

\subsection*{Your turn}
\begin{itemize}
\raggedright
  \item \emph{Say it:} \emph{tobidasu}
  \item \emph{Say it:} \emph{tobidasu}, once more
  \item \emph{Say it:} \emph{kodomo ga michi ni tobidashimasu} — A child rushes out into the street.
  \item \emph{Recall:} say the Japanese for to be crowded, then the Japanese for to cut across, then say \emph{tobidasu} again
\end{itemize}

\begin{tcolorbox}[breakable,colback=teal!4,colframe=teal!35!black,title={Before you move on}]
What does \ja{とびだす} mean? (\textquotedblleft{}To rush out; to jump out\textquotedblright{}.) Say it once more.
\end{tcolorbox}
\section[\texorpdfstring{\ja{ぶつかる} (butsukaru)}{ぶつかる (butsukaru)}]{\ja{ぶつかる} (butsukaru) — to bump into; to collide with}

\label{lesson:JA-C201-butsukaru}



\begin{quote}
Before the new one: say the Japanese for to cut across, then the Japanese for to rush out.
\end{quote}

\subsection*{You'll want to know: \ja{ぶつかる}}
\textbf{\ja{ぶつかる}} — \emph{butsukaru} — \textquotedblleft{}to bump into; to collide with\textquotedblright{}.

\begin{quote}
\textbf{\ja{じてんしゃが} \ja{くるまに} \ja{ぶつかりました。}} — \emph{jitensha ga kuruma ni butsukarimashita} — \textquotedblleft{}A bicycle bumped into a car.\textquotedblright{}
\end{quote}

Polite \textbf{\ja{ぶつかります}}, \emph{butsukarimasu}; a godan verb.

\begin{cousinweb}
\textbf{\ja{ぶつかる}} (\emph{butsukaru}) is a native Japanese word (wago), inherited rather than borrowed and usually written in kana (JMdict); its older origin is not settled.
\end{cousinweb}

\begin{grammarlens}[title={What you've built}]
One more word for everyday life at A2.
\end{grammarlens}

\subsection*{Your turn}
\begin{itemize}
\raggedright
  \item \emph{Say it:} \emph{butsukaru}
  \item \emph{Say it:} \emph{butsukaru}, once more
  \item \emph{Say it:} \emph{jitensha ga kuruma ni butsukarimashita} — A bicycle bumped into a car.
  \item \emph{Recall:} say the Japanese for to cut across, then the Japanese for to rush out, then say \emph{butsukaru} again
\end{itemize}

\begin{tcolorbox}[breakable,colback=teal!4,colframe=teal!35!black,title={Before you move on}]
What does \ja{ぶつかる} mean? (\textquotedblleft{}To bump into; to collide with\textquotedblright{}.) Say it once more.
\end{tcolorbox}
